\documentclass{article}
\usepackage{fontenc}
\usepackage[ngerman]{babel}
\usepackage{amsmath,amssymb,amsthm}
\usepackage{hyperref}

\title{Eine strukturelle Analyse großer Abweichungen in der ABC-Vermutung}
\author{Thomas Hoffbauer}
\date{}

\newtheorem{theorem}{Theorem}
\newtheorem{proposition}{Proposition}
\newtheorem{remark}{Bemerkung}

\begin{document}
	\maketitle
	
	\begin{abstract}
		Wir untersuchen die Größe
		\[
		\Delta(a,b,c) := \log c - \log \mathrm{rad}(abc)
		\]
		für Lösungen der Gleichung $a+b=c$ mit $\gcd(a,b)=1$.
		Numerische und strukturelle Argumente legen nahe, dass große Werte von $\Delta$
		nur bei stark konzentrierten Primfaktorstrukturen auftreten.
	\end{abstract}
	
	\section{Grundlagen}
	
	Für $a+b=c$ gilt:
	\[
	\Delta = \sum_{p \mid c} (\gamma_p - 1)\log p.
	\]
	
	\section{Kontrolle der mehrfachen Primfaktoren}
	
	\begin{proposition}
		Für $\Delta(a,b,c) \ge \delta > 0$ gilt:
		\[
		|S_2| \le \frac{\delta}{\log 2},
		\]
		wobei $S_2 = \{p : \gamma_p \ge 2\}$.
	\end{proposition}
	
	\section{Schranke für quadratfreie Anteile}
	
	\begin{proposition}
		Sei $S_1 = \{p : \gamma_p = 1\}$. Dann gilt:
		\[
		|S_1| \lesssim \frac{\log c}{\log\log c}.
		\]
	\end{proposition}
	
	\section{Heuristische Schranke in $\Delta$}
	
	\begin{proposition}
		Heuristisch gilt:
		\[
		\omega(abc) \lesssim \log \Delta.
		\]
	\end{proposition}
	
	\section{Strukturelles Prinzip}
	
	\begin{theorem}[Strukturprinzip]
		Für $\Delta(a,b,c) \ge \delta > 0$ gilt:
		
		\begin{itemize}
			\item Die Anzahl mehrfacher Primfaktoren ist beschränkt,
			\item die Gesamtanzahl der Primfaktoren wächst höchstens langsam,
			\item die Primfaktorstruktur von $c$ ist stark konzentriert.
		\end{itemize}
	\end{theorem}
	
	\section{Numerische Evidenz}
	
	Berechnungen bis $c \le 2\cdot 10^6$ zeigen:
	
	\begin{itemize}
		\item Große $\Delta$ treten nur bei kleinen $\omega(abc)$ auf,
		\item maximale Werte entstehen bei Potenzstrukturen,
		\item keine Gegenbeispiele mit breiter Primstruktur wurden gefunden.
	\end{itemize}
	
	\section{Interpretation}
	
	Große Abweichungen entsprechen einer Reduktion der effektiven Dimension
	im Raum der Primfaktoren.
	
\end{document}